% =====================================================================
%  Chapter 02 — The three routes
% =====================================================================
\chapter{Three Routes to the Prize}
\label{ch:routes}

The founding thesis of Chapter~\ref{ch:origin} does not by itself say which
system to work in. Temperature is the master variable, but a proof has to be
carried out in some specific set of equations, and the choice of that set is
the choice between being physically natural and being \emph{admissible} for
the Prize. The programme's specification identifies three, ranks them, and
states the limit that carries each back to the Prize equations.

\section{The three systems side by side}
\label{sec:three-systems}

\begin{table}[ht]
\centering
\small
\begin{tabular}{@{}l p{0.27\textwidth} c p{0.22\textwidth} p{0.17\textwidth}@{}}
\toprule
Route & System & Prize geometry? & Final limit & Rank \\
\midrule
\textbf{1} & Incompressible NS with \(\mu(T)\), plus a scalar transport
  equation for \(T\)
  & yes, from the start
  & \(\mu(T)\to\nu\): uniform temperature, thermal relaxation
  & primary \\[4pt]
\textbf{2} & The \emph{exact} Prize equations, plus an auxiliary scalar
  \(\theta\) slaved to the velocity
  & yes, exactly
  & \(\varepsilon\to0\), elementary; and at \(\varepsilon=0\) nothing to take
  & backup \\[4pt]
\textbf{3} & Compressible Navier--Stokes--Fourier
  & no: \(\rho\) is variable
  & \(\mathrm{Ma}\to0\), hard
  & completeness \\
\bottomrule
\end{tabular}
\caption[The three routes]{The three routes, as specified. The ranking is the
programme's own, from its architecture document. It is worth noting that the
route ranked ``backup'' is the one every result in Parts~III--VI is stated
in.}
\label{tab:three-routes}
\end{table}

\section{Route~1: incompressible with temperature-dependent viscosity}
\label{sec:route1}

Route~1 was written down in Chapter~\ref{ch:origin} as
\eqref{eq:route1-u}--\eqref{eq:route1-T}. Its virtue is that it stays in Prize
geometry throughout: \(\rho\) is constant and \(\nabla\cdot u=0\) is enforced
exactly, not approximately, by a spectral Leray projection at every step.
Nothing in Route~1 has to be argued away at the end except the temperature
dependence of \(\mu\) itself.

Its defect is that it is a genuinely different system. Global strong
solvability of the incompressible Navier--Stokes equations with
temperature-dependent viscosity, for large data in three dimensions, is
\emph{open in the literature}; the programme's own list of priority questions
names it first. Route~1 therefore does not reduce the Prize to something
easier. It substitutes a hard open problem for a hard open problem, in
exchange for a mechanism --- the scalar maximum principle --- that the Prize
equations do not obviously offer.

The limit \(\mu(T)\to\nu\) is nevertheless simple in a way that matters. It is
a limit of \emph{coefficients}, not of geometry. Parametrise the initial
temperature as \(T_0=\bar T+\delta\,g(x)\) with \(\norm{g}_{L^\infty}=1\); as
\(\delta\to0\) the temperature field becomes uniform and \(\mu(T)\to\mu(\bar
T)=\nu\). What the programme measured is whether anything degenerates in that
limit, and it does not: Chapter~\ref{ch:layer3} reports
\(\abs{\varepsilon(\delta)}\propto\delta^{0.980}\) with \(R^2>0.99\) across
three decades of \(\delta\), an entirely regular, first-order approach.

\section{Route~2: the exact Prize equations plus a diagnostic scalar}
\label{sec:route2}

Route~2 is the programme's answer to an objection, and understanding the
objection is the fastest way to understand the route.

\subsection{Tao's objection}

Any regularity argument that proceeds by modifying the equations --- adding a
regularising term, changing the viscosity law, working in a different
symmetry class --- owes an account of the limit back to the unmodified
problem. The objection, associated with Terence Tao's writing on the
supercriticality barrier, is that this account is where the difficulty
usually reappears: the modified problem is regular \emph{because} of the
modification, the constants degenerate as the modification is removed, and
nothing has been learned about the original equations. Route~1 is exposed to
exactly this: its regularity mechanism is \(A(T)>0\), and \(A\) is a property
of the modified system.

\subsection{The construction}

Route~2 evades the objection by not modifying anything. Take a solution
\(u\) of the exact Prize equations \eqref{eq:NS-prize} and \emph{define}
\begin{equation}\label{eq:theta-def}
  \partial_t\theta + u\cdot\nabla\theta
    = \nu\Delta\theta + \nu\abs{\nabla u}^2,
  \qquad \theta(\cdot,0)=0 .
\end{equation}
At \(\varepsilon=0\) this scalar does not appear anywhere in
\eqref{eq:NS-prize}. As the proof document states at line~152, it is
\emph{purely diagnostic}: its equation is an \emph{identity satisfied by any
solution}, not an additional equation imposed on one. The velocity field is
completely unaffected by \(\theta\)'s existence.

The optional regularisation --- an effective viscosity
\(\mu_{\mathrm{eff}}=\nu+\varepsilon f(\theta)\) with \(f(0)=0\), \(f\ge0\)
--- exists only so that a limit \(\varepsilon\to0\) can be discussed. If the
argument runs at \(\varepsilon=0\), the limit is vacuous: there is nothing to
take.

\begin{obsv}[The numerics say the limit is vacuous]
\label{obs:eps-vacuous}
The programme tested this directly. In \Sn{05}, the 2D \(\varepsilon\)-sweep
\texttt{EXP-L2-R2-001} through \texttt{-004} at
\(\varepsilon\in\{1,0.1,0.01,0.001\}\) returned
\(\theta_{\max}=2.6045\times10^{-4}\) --- \emph{identical to five significant
figures across all four}, because the source \(\nu\abs{\nabla u}^2\) does not
depend on \(\varepsilon\), and the vorticity evolution is byte-identical
because \(\varepsilon\) never enters it. In \Sn{10} and \Sn{17} the same
comparison was run in 3D including \(\varepsilon=0\) exactly: the regularity
exponent \(\delta_{\max}=1.50\) for Taylor--Green and \(0.50\) for the shear
layer, unchanged at every \(\varepsilon\) including zero
(\texttt{EXP-L3-R2-TG-128}, \texttt{EXP-L3-R2-SH-128}, both at
\texttt{eps\_param}\(=0\), both \textsf{PASS}).
\end{obsv}

That is the whole content of the answer to Tao's objection. Route~2's
regularity statement is a statement about solutions of the exact Prize
equations; the auxiliary scalar is a lens, not an ingredient. Whatever the
route proves, it proves about \eqref{eq:NS-prize}.

\subsection{What Route~2 has to deliver}

The price of working in the exact equations is that no free mechanism comes
with them. Route~2 has to earn its regularity, and the whole of Part~III is
the story of the attempt. The target is a single quantitative statement, which
the proof document names \emph{Claim~A} (\texttt{def:claimA}, line~163):
\begin{equation}\label{eq:claimA}
  \norm{\nu\abs{\nabla u}^2}_{L^{1+\delta}(Q_T)} \le C(\nu,E_0,T)<\infty
  \qquad\text{for some }\delta>0 .
\end{equation}
The energy inequality gives \eqref{eq:claimA} with \(\delta=0\) and nothing
more. Any \(\delta>0\), however small, is enough: the Route~C bootstrap of
\texttt{thm:routeC} (line~508) then drives \(\delta_n\to\infty\), and
Prodi--Serrin gives smoothness. That reduction is unconditional and is one of
the very few things in this book that has never been in doubt.

\section{Route~3: compressible Navier--Stokes--Fourier}
\label{sec:route3}

Route~3 is the full compressible system with a genuine energy equation, and
it is where the T-first thesis is most at home: temperature is a real
thermodynamic variable, \(\rho\), \(p\), \(\mu\), \(k\) are all functions of
the state, and the second law is not an assumption but a structural
constraint on the system.

It was demoted for one reason, and the reason is decisive: the density is not
constant. The Prize problem is posed for the incompressible equations with
\(\nabla\cdot u=0\); a compressible solution reaches them only through the
low-Mach-number limit \(\mathrm{Ma}\to0\). That limit is a genuine research
area in its own right --- it involves acoustic waves whose speed diverges,
initial layers, and delicate uniform estimates --- and it is exactly the kind
of place where, per Tao's objection, the difficulty could reappear entirely
intact.

Route~3 therefore became an \emph{independent completeness argument}: evidence
that the T-first mechanism is not an artefact of the incompressible
idealisation, and nothing more. Two locked baseline experiment suites,
\texttt{verified/tfirst\_verify.py} (ideal gas, five experiments) and
\texttt{verified/sco2\_verify.py} (supercritical CO\textsubscript{2} via
\textsc{CoolProp}, seven claims), belong to this line and pass. They are not
cited in support of any regularity claim anywhere in this book.

\section{How the three routes relate}
\label{sec:routes-relate}

\begin{figure}[ht]
\centering
\begin{tikzpicture}[
  >={Latex[length=2.2mm]},
  sys/.style={draw, thick, rounded corners=3pt, align=center, inner sep=5pt,
              font=\small, minimum width=42mm, minimum height=13mm},
  prize/.style={draw, very thick, rounded corners=3pt, align=center,
                inner sep=6pt, font=\small\bfseries, fill=panelgrey,
                minimum width=44mm},
  lim/.style={font=\scriptsize, align=center, inner sep=2pt}
]
\node[sys] (r1) at (0, 3.0)
  {\textbf{Route 1}\\[1pt]
   \(\nabla\!\cdot\!(\mu(T)\nabla u)\), scalar \(T\)\\[1pt]
   \(\nabla\!\cdot\!u=0\), \(\rho\) const};
\node[sys] (r2) at (0, 0.4)
  {\textbf{Route 2}\\[1pt]
   exact \eqref{eq:NS-prize} \(+\) scalar \(\theta\)\\[1pt]
   \(\mu_{\mathrm{eff}}=\nu+\varepsilon f(\theta)\)};
\node[sys] (r3) at (0,-2.3)
  {\textbf{Route 3}\\[1pt]
   compressible NSF\\[1pt]
   \(\rho\) variable};

\node[prize] (P) at (8.4, 0.4)
  {The Prize equations\\
   \(\partial_tu+(u\!\cdot\!\nabla)u=-\nabla p+\nu\Delta u\)};

\draw[->, thick] (r1.east) to[out=0, in=140]
  node[lim, above, pos=0.42, sloped] {\(\mu(T)\to\nu\)}
  node[lim, below, pos=0.42, sloped] {simple} (P.north west);

\draw[->, very thick] (r2.east) --
  node[lim, above] {\(\varepsilon\to0\); at \(\varepsilon=0\), \emph{already there}}
  (P.west);

\draw[->, thick, dashed] (r3.east) to[out=0, in=-140]
  node[lim, above, pos=0.42, sloped] {\(\mathrm{Ma}\to0\)}
  node[lim, below, pos=0.42, sloped] {hard} (P.south west);

\node[font=\scriptsize\itshape, text=rulegrey, right=2mm of r3, yshift=-9mm,
      align=left, text width=52mm]
  {demoted to an independent completeness argument: the limit is where the
   difficulty could reappear};
\end{tikzpicture}
\caption[The three routes and their limits]{The three routes and the limits
that carry each back to the Prize equations. Route~2's arrow is drawn heavy
because it is not really a limit at all: at \(\varepsilon=0\) the system
\emph{is} the Prize system, and the auxiliary scalar is a diagnostic
identity.}
\label{fig:routes}
\end{figure}

The relation between them, stated once and for the record: Route~1 buys a
mechanism at the cost of changing the equations; Route~2 keeps the equations
and must find its mechanism; Route~3 keeps the physics and must pay for it in
the limit. Every route was built, tested and instrumented
(Chapters~\ref{ch:layer12}--\ref{ch:layer4}). Every analytical result the
programme produced after \Sn{18} is a Route~2 result.

\begin{record}
\textbf{The correction the reader should carry forward.} The programme's own
architecture calls Route~1 primary and Route~2 backup. That ranking was
never formally revised, and it is still what the specification says. It is
also, by the evidence of the proof document, wrong: from \Sn{18} to \Sn{47}
the entire analytical effort is Route~2, working inside the exact Prize
equations at constant viscosity. Route~1 remains a fully built and tested
numerical programme with no analytical theory attached to it.
\end{record}
